The Kerr-Schild coordinates are labels of the flat background and $g = \eta - 2H\,k \otimes k$ is the derived metric of $\hb$ (Proposition \claimref{interior.unipotency-across-kerr-family}; Fact \claimref{synthesis.imported-gtg-dictionary}(i)). The physical vector over a background vector $b$ is $\theta(b) := \hb^{-1}(b) = b - H(b \cdot k)k$, since $k$ is null; $\hb$ and $\theta$ are symmetric for the flat product. Physical vectors are flat vectors, and $\Gamma(A; B, C) = \bigl(\hb(F_A) \cdot \nabla F_B + \Omega(\hb(F_A)) \cdot F_B\bigr) \cdot F_C$, by the covariant derivative of Definition \claimref{obj.synthesis.gtg-flat-spacetime-gauge-fields}(b) with $B \times u = B \cdot u$ for a bivector $B$ and a vector $u$. We use $((p \wedge u) \cdot v) \cdot x = (p \cdot x)(u \cdot v) - (p \cdot v)(u \cdot x)$, so for $\Omega(b) = \dot\nabla \wedge U_b$ with $U_b = \dot H(b \cdot \dot k)\dot k$,
\[
((\Omega(b) \cdot v) \cdot x) = (x \cdot \nabla U_b) \cdot v - (v \cdot \nabla U_b) \cdot x . \tag{$*$}
\]
Fact \claimref{interior.imported-kerr-schild-form-kerr-gauge} gives $k^\mu = (1, (rx + ay)/(r^2 + a^2), (ry - ax)/(r^2 + a^2), Z/r)$, null, with spatial part $v$ constant along the straight lines it is tangent to, so $k \cdot \nabla k = 0$; $k$ and $v$ are stationary.

\emph{(1) Exact connection.} For spinless matter $\Omega$ is fixed by vanishing torsion (Definition \claimref{obj.synthesis.gtg-flat-spacetime-gauge-fields}(f)); for coordinate vector fields this is the symmetry in $b, c$ of $b \cdot \nabla\theta(c) + \Omega(b) \cdot \theta(c)$, that is
\[
b \cdot \nabla\theta(c) - c \cdot \nabla\theta(b) + \Omega(b) \cdot \theta(c) - \Omega(c) \cdot \theta(b) = 0
\]
for constant $b, c$, and a bivector-valued solution is unique (the Levi-Civita connection). Put $h(b) = H(b \cdot k)k$, symmetric, so $\theta = 1 - h$ and $\Omega(b) = \dot\nabla \wedge h(b)$ with $b$ fixed. By $(*)$, $\Omega(b) \cdot c = \dot\nabla(\dot h(b) \cdot c) - c \cdot \nabla h(b)$, and the symmetry of $h$ gives $\Omega(b) \cdot c - \Omega(c) \cdot b = b \cdot \nabla h(c) - c \cdot \nabla h(b)$, which cancels the terms linear in $h$. The remainder is $-\Omega(b) \cdot h(c) + \Omega(c) \cdot h(b) = -H(c \cdot k)\,\Omega(b) \cdot k + H(b \cdot k)\,\Omega(c) \cdot k$. Now $\Omega(b) \cdot k = \dot\nabla(\dot h(b) \cdot k) - k \cdot \nabla h(b)$; the first term is a derivative of $H(b \cdot k)(k \cdot k)$ with $k$ differentiated only in $h$, which vanishes because $k \cdot k = 0$ implies $(u \cdot \nabla k) \cdot k = 0$; and $k \cdot \nabla h(b) = (k \cdot \nabla H)(b \cdot k)k$ because $k \cdot \nabla k = 0$. So $\Omega(b) \cdot k = -(b \cdot k)(k \cdot \nabla H)\,k$, the remainder is $H(k \cdot \nabla H)\bigl((c \cdot k)(b \cdot k) - (b \cdot k)(c \cdot k)\bigr)k = 0$, and the torsion equation holds exactly. Only $k \cdot k = 0$ and $k \cdot \nabla k = 0$ were used.

\emph{(2) Frame components.} By $(*)$ and (1), with $\alpha_A := F_A \cdot k$ and $F_A$ frozen, $((\Omega(\hb F_A) \cdot F_B) \cdot F_C) = (F_C \cdot \nabla U) \cdot F_B - (F_B \cdot \nabla U) \cdot F_C$, $U = H\alpha_A k$, and $(u \cdot \nabla U) \cdot F_B = (u \cdot \nabla H)\alpha_A(k \cdot F_B) + H\,(F_A \cdot (u \cdot \nabla k))(k \cdot F_B) + H\alpha_A\,(u \cdot \nabla k) \cdot F_B$. Since $v \cdot v$ is constant and $e_0$ is constant, $(u \cdot \nabla v) \cdot v = 0 = (u \cdot \nabla v) \cdot e_0$ for every $u$, so $(u \cdot \nabla k) \cdot k = (u \cdot \nabla k) \cdot \ell = 0$ and $\ell \cdot (u \cdot \nabla k) = 0$. On stationary functions $k \cdot \nabla = v^i\partial_i$ and $\ell \cdot \nabla = -\tfrac12 v^i\partial_i$; $k \cdot \nabla v = 0$ gives $\ell \cdot \nabla v = 0$, and since $e_1, e_2$ are functions of $v$, $k \cdot \nabla$ and $\ell \cdot \nabla$ annihilate the whole frame. For $A = k$: $\hb(k) = k$, $\alpha_k = 0$ and $F_A \cdot (u \cdot \nabla k) = 0$, so $\Omega(k) = 0$ and $\Gamma(k; B, C) = (k \cdot \nabla F_B) \cdot F_C = 0$. For $A = e_i$: $\hb(e_i) = e_i$, $\alpha = 0$, and the connection term is $H\,(e_i \cdot (F_C \cdot \nabla v))(k \cdot F_B) - (B \leftrightarrow C)$, nonzero only when $\ell \in \{F_B, F_C\}$; so $\Gamma(e_i; k, \ell) = (e_i \cdot \nabla k) \cdot \ell - H\,\bigl(e_i \cdot (k \cdot \nabla v)\bigr)(k \cdot \ell) = 0$, and $\Gamma(e_i; \ell, e_j) = (e_i \cdot \nabla \ell) \cdot e_j + H\,e_i \cdot (e_j \cdot \nabla v)$. For $A = \ell$: $\hb(\ell) = \ell + Hk$, $\alpha = 1$, and
\[
\Gamma(\ell; B, C) = (\ell \cdot \nabla F_B) \cdot F_C + (F_C \cdot \nabla H)(k \cdot F_B) - (F_B \cdot \nabla H)(k \cdot F_C) + H\bigl((F_C \cdot \nabla v) \cdot F_B - (F_B \cdot \nabla v) \cdot F_C\bigr),
\]
which gives $\Gamma(\ell; k, e_i) = 0$, $\Gamma(\ell; \ell, e_i) = e_i \cdot \nabla H$, $\Gamma(\ell; k, \ell) = -k \cdot \nabla H$ and $\Gamma(\ell; e_1, e_2) = H\bigl((e_2 \cdot \nabla v) \cdot e_1 - (e_1 \cdot \nabla v) \cdot e_2\bigr)$. From $z^2 = w$ (Theorem \claimref{synthesis.kerr-ring-double-cover-data}(1)), $\partial_x z = x/z$, $\partial_y z = y/z$, $\partial_Z z = (Z - Ia)/z$, so $H = M\operatorname{Re}(1/z)$ has $\partial H = -M\operatorname{Re}((x, y, Z - Ia)z^{-3})$; and with $v = \bigl(rx + ay,\ ry - ax,\ (r^2 + a^2)\cos\theta\bigr)/(r^2 + a^2)$ and $x^2 + y^2 = (r^2 + a^2)\sin^2\theta$, $v^i\partial_i z = \bigl(r\sin^2\theta + \cos\theta\,(r\cos\theta - Ia)\bigr)/z = 1$.

\emph{(3) Orders.} The weights $(\tfrac12, -\tfrac12)$ of Theorem \claimref{synthesis.kerr-ring-blow-up-forced-cover-fixed-order}(2) are exponents on cometric components; on the frame vectors they are $F_K = \rho^{-1/2}k$, $F_N = \rho^{1/2}\ell$, the same rescaling. On the cover $r = \operatorname{Re} z$ and $a\cos\theta = -\operatorname{Im} z$, and $v$, $e_1$, $e_2$, $x$, $y$, $Z$ are smooth functions of $(\varphi, z)$ up to the ring (Theorem \claimref{synthesis.kerr-ring-blow-up-forced-cover-fixed-order}(1)); by the formulas for $\partial z$ their first Cartesian derivatives are $O(\rho^{-1})$, uniformly in $\varphi$ and time. Hence $\nabla v, \nabla e_i = O(\rho^{-1})$, $H = O(\rho^{-1})$, $\nabla H = O(\rho^{-3})$, $k \cdot \nabla H = M\operatorname{Re}(-z^{-2}) = O(\rho^{-2})$, $\nabla\log\rho = \operatorname{Re}(\nabla z/z) = O(\rho^{-2})$, $k \cdot \nabla\log\rho = \operatorname{Re}(1/z) = O(\rho^{-1})$ and $\ell \cdot \nabla\log\rho = -\tfrac12\operatorname{Re}(1/z)$. The remaining unweighted components are $\Gamma(e_i; k, e_j) = (e_i \cdot \nabla v) \cdot e_j = O(\rho^{-1})$, $\Gamma(e_i; \ell, e_j) = -\tfrac12(e_i \cdot \nabla v) \cdot e_j + H\,e_i \cdot (e_j \cdot \nabla v) = O(\rho^{-2})$, $\Gamma(e_i; e_1, e_2) = (e_i \cdot \nabla e_1) \cdot e_2 = O(\rho^{-1})$ and $\Gamma(\ell; e_1, e_2) = O(\rho^{-2})$. Weighting multiplies $\Gamma(A; B, C)$ by $\rho^{w_A + w_B + w_C}$ with $w = (-\tfrac12, \tfrac12, 0, 0)$, and adds, from $D(\rho^{\mp 1/2}u) = \rho^{\mp1/2}(Du \mp \tfrac12 (d\log\rho) u)$, the term $-\tfrac12\rho^{w_A}\,\hb(F_A) \cdot \nabla\log\rho$ to $\Gamma(A; F_K, F_N)$ (and its negative to $\Gamma(A; F_N, F_K)$). The weighted orders are therefore: $\Gamma(F_K; F_K, F_N) = -\tfrac12\rho^{-1/2}k \cdot \nabla\log\rho = O(\rho^{-3/2})$ and all other components along $F_K$ zero; $\Gamma(F_N; F_K, F_N) = \rho^{1/2}\bigl(-k \cdot \nabla H - \tfrac12(\ell + Hk) \cdot \nabla\log\rho\bigr) = O(\rho^{-3/2})$, $\Gamma(F_N; F_K, e_i) = 0$, $\Gamma(F_N; F_N, e_i) = \rho\,e_i \cdot \nabla H = O(\rho^{-2})$, $\Gamma(F_N; e_1, e_2) = O(\rho^{-3/2})$; $\Gamma(e_i; F_K, F_N) = -\tfrac12 e_i \cdot \nabla\log\rho = O(\rho^{-2})$, $\Gamma(e_i; F_K, e_j) = O(\rho^{-3/2})$, $\Gamma(e_i; F_N, e_j) = O(\rho^{-3/2})$, $\Gamma(e_i; e_1, e_2) = O(\rho^{-1})$. By (4) below the leading coefficient of $\Gamma(F_N; F_N, e_i)$ has modulus $M/2 \neq 0$ in every direction $\psi$, so the order $-2$ is attained on the tangential direction $F_N$; a normalisation $\rho^m$ keeps every component bounded only for $m \geq 2$ and leaves a nonzero tangential limit only for $m = 2$; $d = |w|/(2a)$ is $\rho^2/(2a)$, the flat distance to the ring to leading order since $dw = 2a\,d\varpi - 2Ia\,dZ$ there.

\emph{(4) The limit.} Multiplying by $d$, every component listed as $O(\rho^{-3/2})$ or smaller tends to zero like $O(\rho^{1/2})$, which gives $\widehat\Omega(F_K) \to 0$. For the transverse vectors, $d\,\Gamma(e; F_K, F_N) = -(e \cdot \nabla|w|)/(8a)$ is linear in $e$; it vanishes on $\hat e_\psi$ by definition, and on $\nu$, since $\nabla|w| \to 2a\,\hat n$ with $\hat n$ the outward transverse normal at the ring, it tends to $-\tfrac14$. As $((F_K \wedge F_N) \cdot F_K) \cdot F_N = 1$, this is $\widehat\Omega(\nu) \to -\tfrac14 F_K \wedge F_N$ and $\widehat\Omega(\hat e_\psi) \to 0$. Along $F_N$ the surviving components are $d\rho\,e_i \cdot \nabla H = (\rho^3/2a)\,\partial_{e_i}H$. At the ring point $(x, y, Z) = a\hat\varpi$ and $e_1 \to \hat z$, $e_2 \to \hat\varpi$ up to $O(\rho)$, so $\hat\varpi \cdot (x, y, Z - Ia) \to a$ and $\hat z \cdot (x, y, Z - Ia) \to -Ia$, and
\[
\tfrac{\rho^3}{2a}\partial_{\hat\varpi}H \to -\tfrac{M}{2}\operatorname{Re}(\rho^3z^{-3}), \qquad \tfrac{\rho^3}{2a}\partial_{\hat z}H \to -\tfrac{M}{2}\operatorname{Re}(-I\rho^3z^{-3}) = -\tfrac{M}{2}\operatorname{Im}(\rho^3z^{-3}),
\]
so $L_\varpi + IL_Z = -\tfrac{M}{2}\rho^3z^{-3} = -\tfrac{M}{2}e^{-3I\psi} = \tfrac12\rho^3\mathfrak M$ (Proposition \claimref{interior.complex-scale-inverse-cube-killing-spinor}; Theorem \claimref{synthesis.kerr-ring-double-cover-data}(4)). A bivector with $((X \cdot F_N) \cdot e_i) = L_i$ and all other frame components zero is $X = \sum_i L_i\,F_K \wedge e_i$, since $((F_K \wedge e) \cdot F_N) \cdot e' = -(e \cdot e')$. Hence $X(\psi) = F_K \wedge (L_\varpi\hat\varpi + L_Z\hat z)$. With $K = e_0 + \hat s$: $K e_0\hat s = K$, and $\hat\varpi$ commutes with $e_0\hat s$, so $(K \wedge \hat\varpi)I = K e_0 \hat s\,\hat\varpi\hat\varpi\hat z = -K \wedge \hat z$; so $X = (F_K \wedge \hat\varpi)(L_\varpi - IL_Z) = -\tfrac{M}{2}(F_K \wedge \hat\varpi)e^{3I\psi}$. All estimates are uniform in $\psi$, $\varphi$ and time.

\emph{(5) Algebra, duality, winding.} The values of $\widehat\Omega$ on tangential directions tend to elements of $\mathfrak n$; at $\psi = 0$ and $\psi = \pi/6$ they are $F_K \wedge \hat\varpi$ and $F_K \wedge \hat z$ up to nonzero factors, so they span $\mathfrak n$. Adding the normal direction adds the value $-\tfrac14 F_K \wedge F_N$ of (4); since $(F_K \wedge F_N) \times (F_K \wedge e) = F_K \wedge e$ for $e \perp F_K, F_N$ (the boost acts on $\mathfrak n$ with weight one), all limiting values span $\mathfrak n + \R\,F_K \wedge F_N$, solvable with derived algebra $\mathfrak n$, and not nilpotent. For $B = K \wedge e$ with $e \perp K$ spacelike: $B \times K = 0$; $(K \wedge e) \times (K \wedge e') = 0$ for $e, e' \perp K$ (both products reduce to multiples of $K^2 = 0$; equivalently $R_u = 1 + Ku$ with $u \in \mathrm{span}(\hat\varpi, \hat z)$ form the abelian group $(U, +)$ of Proposition \claimref{transport.group-null-rotations}, whose Lie algebra is $\mathfrak n$); and $B \times (B \times (B \times u)) = 0$ while $B \times (B \times N) = -Q(e)\,K \neq 0$ for $N$ with $K \cdot N = 1$, $N \perp e$: null rotations about $K$, nilpotent of order three. The pseudoscalar commutes with bivectors and $(K \wedge \hat\varpi)I = -K \wedge \hat z$, $(K \wedge \hat z)I = K \wedge \hat\varpi$, so $\mathfrak n I = \mathfrak n$ and $e^{I\alpha}$ acts on $\mathfrak n$ as rotation by $\alpha$. The modulus of $X$ in the basis $(F_K \wedge \hat\varpi, F_K \wedge \hat z)$ is $M/2$, and its phase is $3\psi$ plus a constant, so it has degree three round the circle of $\psi$; in the complex coordinate $L_\varpi + IL_Z = -\tfrac{M}{2}e^{-3I\psi}$ the degree is $-3$, the sign recorded for $\rho^3\mathfrak M$ in Theorem \claimref{synthesis.kerr-ring-double-cover-data}(4). A rotation of $(e_1, e_2)$ by an angle $\beta$, and a boost $(k, \ell) \mapsto (\lambda k, \lambda^{-1}\ell)$, both continuous on the cover up to the ring, change $X$ by the phase $e^{I\beta}$ and the positive factor $\lambda^{\pm1}$ in the limit; on the cover the circle of $\psi$ bounds a disc on which $\beta$ is continuous, so $\beta$ has degree zero round it and the degree of $X$ is unchanged. The deck transformation is $\psi \mapsto \psi + \pi$ at a fixed base point (Theorem \claimref{synthesis.kerr-ring-double-cover-data}(1)), under which $e^{3I\psi}$ changes sign, and $k$, hence the frame, agrees on the two sheets at the ring; so $X$ is odd and downstairs two-valued with holonomy $-1$. For the tractor gauge field $\Omega_T(a) = \Omega(a) - \nInf \wedge \hb^{-1}(a) + \nO \wedge \Sch(\hb^{-1}(a))$ (Theorem \claimref{synthesis.gtg-tractor-gauge-field-reproduces-tractor-connection}): Kerr is vacuum with $\Lambda = 0$, so $\Sch = 0$; on a frame vector $F_A$ the soldering term $-\nInf \wedge F_A$ has frame components independent of $\rho$, so $d$ times it tends to zero; and $X \in \Lambda^2 W$ commutes with $\nInf$ and $\nO$, which are orthogonal to $W$. So the normalised tractor gauge field has the same limit, and its values annihilate the scale tractor $\nInf$.

\emph{(6) Linking loop.} Near the ring point write the meridian plane as $\C$ with $\hat\varpi \leftrightarrow 1$, $\hat z \leftrightarrow I$ and $P = (\varpi - a) + IZ$. From $w = z^2$ and $dw = 2a\,d\varpi - 2Ia\,dZ$ at the ring, $P = \bar z^2/(2a)$ to leading order, and at fixed azimuth $v = \hat s + v_\perp$ with $v_\perp = \bar z/a$ to first order (from $v = (rx + ay, ry - ax, Z(r^2 + a^2)/r)/(r^2 + a^2)$ at the azimuth of $\hat\varpi$), so $v_\perp$ is a holomorphic function of $P$ with $dv_\perp/dP = 1/\bar z$. On the loop $dz/d\psi = Iz$, so $T_\perp = dP/d\psi = -I\bar z^2/a$ and $T$ has no $\hat s$ component; $|T| = \rho^2/a$. By the formula for $\Gamma(\ell; B, C)$ in (2) with $\hb(\ell)$ replaced by $T$ ($\alpha = T \cdot k$ no longer zero), $(D_T \ell) \cdot e = (e \cdot \nabla H)(T \cdot k) + H\,T \cdot (e \cdot \nabla v) + O(\rho)$, the term $(T \cdot \nabla\ell) \cdot e = O(|T||\nabla v|) = O(\rho)$ and all other components being $O(\rho)$ by the orders in (3) times $|T|$. For spatial vectors the flat product is minus the Euclidean one, so $T \cdot k = -\operatorname{Re}(T_\perp\bar v_\perp) = \rho^2\operatorname{Im}z/a^2$ and $H\,T \cdot (e \cdot \nabla v) = -H\operatorname{Re}(\epsilon\,\overline{T_\perp/z})$ for $e \leftrightarrow \epsilon$, while $e \cdot \nabla H = \operatorname{Re}(\epsilon\,\overline{\nabla H})$ with $\nabla H = -aMz^{-3}$ by (4). The component along $e$ is therefore $\operatorname{Re}(\epsilon\bar Q)$ with
\[
Q = (T \cdot k)\,\nabla H - H\,T_\perp/z = -\tfrac{M}{a}\sin\psi\,e^{-3I\psi} + I\tfrac{M}{a}\cos\psi\,e^{-3I\psi} = I\tfrac{M}{a}e^{-2I\psi},
\]
so $G_\varpi + IG_Z = Q$, and since $((k \wedge e) \cdot \ell) \cdot e' = -(e \cdot e')$ the generator is $A_0 = k_0 \wedge (G_\varpi\hat\varpi + G_Z\hat z)$. Its values lie in $\mathfrak n$, which is abelian, so the transport of $A_0$ over $[0, \psi]$ is $\exp\int_0^\psi A_0$, and by Gronwall the transport of $A_0 + O(\rho)$ differs from it by $O(\rho)$ over the bounded interval. As $\int_0^{2\pi}e^{-2I\psi}d\psi = \int_0^{\pi}e^{-2I\psi}d\psi = 0$, both transports tend to the identity; a holonomy changes by conjugation under a change of frame at the base point, so its limit, the identity, is frame-independent.

\emph{The checks.} At rational Kerr-Schild points on both sheets ($r > 0$ and $r < 0$), with exact first jets from implicit differentiation of $r(x, y, Z)$: the formula of (1) satisfies the torsion equation, agrees with an independent Koszul solution built from the Lie brackets of $\hb(e_a)$, and gives $\Omega(b) \cdot k = -(b \cdot k)(k \cdot \nabla H)k$; the identities of (2) hold and $\nabla H = -M\operatorname{Re}((x, y, Z - Ia)z^{-3})$. Along rational curves $r = \alpha t$, $q = 1 + \gamma t$ to the ring ($\cos\theta = (1 - q^2)/(1 + q^2)$), every weighted component is an exact rational function of $t$ times a power of $\rho^{1/2}$, with $\rho^2 = (\alpha^2 + a^2\gamma^2)t^2 + O(t^3)$: the top order is $-2$, attained exactly on $\Gamma(F_N; F_N, e_i)$ and the boost along $e_i$ for those generic directions; the leading coefficients $\gamma_\varpi, \gamma_Z$ of $\Gamma(\ell; \ell, e_2)$, $\Gamma(\ell; \ell, e_1)$ satisfy $(\gamma_\varpi + I\gamma_Z)(\alpha + Ia\gamma)^3 = -aM$ exactly, which is $L_\varpi + IL_Z = -\tfrac{M}{2}e^{-3I\psi}$; the radial boost has modulus $\tfrac14$; and on the deck-image curve $(-\alpha, -\gamma)$ the leading coefficients change sign. For the loop tangent $T$ along two curves, every component of $(D_T F_B) \cdot F_C$ is $O(t)$ except $(D_T\ell) \cdot e_i$, whose limits satisfy $(G_\varpi + IG_Z)(\alpha + Ia\gamma)^2 = I\tfrac{M}{a}(\alpha^2 + a^2\gamma^2)$ exactly. In $\Cl(1,3)$ and $\Cl(2,4)$ the plane $\mathfrak n$ is checked abelian, $I$-stable, nilpotent of order three, and commuting with $\nInf$ and $\nO$. As a float cross-check, not part of the evidence, the Ricci contraction of the field strength of (1) computed by finite differences vanishes to $3 \times 10^{-7}$ at a sample point, for $a = 0$ and $a = 1$.
